\section{Limitations}
\label{sec:b-limitations}

\begin{itemize}[leftmargin=1.4em]
\item \textbf{One-step scope.} The controller acts once, at one checkpoint,
  under a frozen continuation.  Deployment value therefore describes a
  controller that intervenes at a checkpoint drawn as the design draws it; a
  stopping-time or every-checkpoint deployment adds selection and coverage terms
  the design does not estimate (\cref{sec:bellman}).  No global guarantee over
  reachable states is claimed.
\item \textbf{Fault prevalence.} Controlled mechanisms establish that a
  difference can exist under specified conditions; they do not establish
  prevalence in natural deployments.  If only controlled faults support the
  result, the paper is framed accordingly, and no-injected-fault results are
  reported separately rather than averaged in.
\item \textbf{Feasibility is a registry claim.} The mask is only as valid as its
  registry and evaluator.  Withholding verifier access synthetically is an
  experimental capacity condition, not evidence of natural unavailability, and
  it is recorded whether it changes available information, cost or baseline
  behaviour.
\item \textbf{Frozen continuation and representation.} Values are relative to
  the registered $\rho$ and horizon.  The registered representation may discard
  decision-relevant history; representation adequacy is a separate estimable
  question and is not assumed by the controller.
\item \textbf{Finite data and multiplicity.} Effective sample sizes at small
  capacities can be much smaller than root counts; per-cell claims are reported
  with their support and abstain when support is insufficient.  A cluster
  bootstrap is asymptotic.
\item \textbf{Substrate and external validity.} Effects are scoped to the
  specified model, harness, continuation and task distribution.  The first study
  identifies no effects for another model, prompt, tool policy or workload; a
  second deployment is a scoped replication, not an automatic broad claim.
\item \textbf{Safety language.} Empirical protection, certified noninferiority
  and individual harm probability are distinct claims with distinct evidence
  requirements (\cref{sec:b-conservative}); none is upgraded by a passing test
  or a LaTeX build.
\item \textbf{Batch versus streaming.} Batch results do not establish streaming
  rate guarantees; the two are reported separately.
\end{itemize}
